\documentclass[twocolumn]{aastex631}

\usepackage{amsmath}
\usepackage{multirow}
\usepackage{natbib}
\usepackage{graphicx} 
\usepackage{aas_macros}

\begin{document}

Understanding the quantum stability of massive gravity theories is crucial for developing consistent extensions of General Relativity. This paper rigorously investigates a specific class of all-order stable massive Degenerate Higher-Order Scalar-Tensor (DHOST) theories, aiming to precisely delineate the mechanisms guaranteeing their classical and quantum health. We perform a formal classification of these theories within the canonical DHOST framework, identifying the minimal structural conditions for their exceptional stability, and then employ an Effective Field Theory (EFT) approach to analyze quantum stability at the one-loop level, focusing on protective symmetries and the graviton mass. Finally, we examine the massless limit of the theory. Our analysis reveals that these theories constitute a fine-tuned subclass of Class N-I DHOST models, characterized by a constant gravitational coupling, a specific inverse kinetic term dependence for the 'beyond Horndeski' function ($F_4 \sim 1/X$), and an internal symmetry relation. The EFT analysis identifies a crucial non-linear Galilean-type shift symmetry in the kinetic sector that functions as a non-renormalization theorem, forbidding the one-loop generation of dangerous ghost-reintroducing or strong-coupling-lowering operators. We demonstrate that the graviton mass term, while explicitly breaking this symmetry, does so 'softly,' ensuring that any radiatively generated symmetry-violating operators are suppressed by powers of the mass scale, thereby preserving technical naturalness and quantum stability. Furthermore, the massless limit is shown to be unique and unambiguous, smoothly recovering a well-defined, ghost-free massless 'beyond Horndeski' theory with complete Stueckelberg field decoupling.

\end{document}
                